\documentclass{amsart}
% For dealing with references we use the comment environment
\usepackage{verbatim}
\newenvironment{reference}{\comment}{\endcomment}
%\newenvironment{reference}{}{}
\newenvironment{slogan}{\comment}{\endcomment}
\newenvironment{history}{\comment}{\endcomment}

% For commutative diagrams we use Xy-pic
\usepackage[all]{xy}

% We use 2cell for 2-commutative diagrams.
\xyoption{2cell}
\UseAllTwocells

% We use multicol for the list of chapters between chapters
\usepackage{multicol}

% This is generally recommended for better output
\usepackage{lmodern}
\usepackage[T1]{fontenc}

% For cross-file-references

% Package for hypertext links:
\usepackage{hyperref}

% For any local file, say "hello.tex" you want to link to please

% Theorem environments.
%
\theoremstyle{plain}
\newtheorem{theorem}[subsection]{Theorem}
\newtheorem{proposition}[subsection]{Proposition}
\newtheorem{lemma}[subsection]{Lemma}

\theoremstyle{definition}
\newtheorem{definition}[subsection]{Definition}
\newtheorem{example}[subsection]{Example}
\newtheorem{exercise}[subsection]{Exercise}
\newtheorem{situation}[subsection]{Situation}

\theoremstyle{remark}
\newtheorem{remark}[subsection]{Remark}
\newtheorem{remarks}[subsection]{Remarks}

\numberwithin{equation}{subsection}

% Macros
%
\def\lim{\mathop{\mathrm{lim}}\nolimits}
\def\colim{\mathop{\mathrm{colim}}\nolimits}
\def\Spec{\mathop{\mathrm{Spec}}}
\def\Hom{\mathop{\mathrm{Hom}}\nolimits}
\def\Ext{\mathop{\mathrm{Ext}}\nolimits}
\def\SheafHom{\mathop{\mathcal{H}\!\mathit{om}}\nolimits}
\def\SheafExt{\mathop{\mathcal{E}\!\mathit{xt}}\nolimits}
\def\Sch{\mathit{Sch}}
\def\Mor{\mathop{\mathrm{Mor}}\nolimits}
\def\Ob{\mathop{\mathrm{Ob}}\nolimits}
\def\Sh{\mathop{\mathit{Sh}}\nolimits}
\def\NL{\mathop{N\!L}\nolimits}
\def\CH{\mathop{\mathrm{CH}}\nolimits}
\def\proetale{{pro\text{-}\acute{e}tale}}
\def\etale{{\acute{e}tale}}
\def\QCoh{\mathit{QCoh}}
\def\Ker{\mathop{\mathrm{Ker}}}
\def\Im{\mathop{\mathrm{Im}}}
\def\Coker{\mathop{\mathrm{Coker}}}
\def\Coim{\mathop{\mathrm{Coim}}}

% Boxtimes
%
\DeclareMathSymbol{\boxtimes}{\mathbin}{AMSa}{"02}

%
% Macros for moduli stacks/spaces
%
\def\QCohstack{\mathcal{QC}\!\mathit{oh}}
\def\Cohstack{\mathcal{C}\!\mathit{oh}}
\def\Spacesstack{\mathcal{S}\!\mathit{paces}}
\def\Quotfunctor{\mathrm{Quot}}
\def\Hilbfunctor{\mathrm{Hilb}}
\def\Curvesstack{\mathcal{C}\!\mathit{urves}}
\def\Polarizedstack{\mathcal{P}\!\mathit{olarized}}
\def\Complexesstack{\mathcal{C}\!\mathit{omplexes}}
% \Pic is the operator that assigns to X its picard group, usage \Pic(X)
% \Picardstack_{X/B} denotes the Picard stack of X over B
% \Picardfunctor_{X/B} denotes the Picard functor of X over B
\def\Pic{\mathop{\mathrm{Pic}}\nolimits}
\def\Picardstack{\mathcal{P}\!\mathit{ic}}
\def\Picardfunctor{\mathrm{Pic}}
\def\Deformationcategory{\mathcal{D}\!\mathit{ef}}

\setlength{\emergencystretch}{3em}
\begin{document}
\title[Stacks Project, Tag 09Y9]{280 --- Finite representable coequalizer presentations of constructible etale sheaves}
\maketitle
\noindent Copyright (C) 2005--2025 Johan de Jong. Stacks Project authors. Source: \href{https://stacks.math.columbia.edu/tag/09Y9}{Tag 09Y9}.

\noindent Editorial difficulty note (not author text). Difficulty H2: A finite stratification alone does not furnish a global presentation by etale representables. The proof must extend stalk elements from the closed stratum to ambient etale neighbourhoods, use quasi-compactness to make finitely many lifts suffice, and inductively combine this with the open strata; the kernel relation then supplies the coequalizer..

\noindent Editorial provenance note (not author text). History: excerpt from revision \texttt{a0444\allowbreak 6e57e\allowbreak c1fbc\allowbreak 252a8\allowbreak 71afc\allowbreak ec775\allowbreak 2fb28\allowbreak 07b14}, etale-cohomology.tex, lines 12342--12409. Reference presentation was adapted; original-author context and core support blocks were retained or added. The mathematical wording is unchanged. Licensed under GNU FDL 1.2 or later; no invariant sections or cover texts. See ../licenses/COPYING. Context blocks reproduce author definitions and supporting results; they are not additional counted problems.

\begin{lemma}
\label{lemma-category-constructible-sets}
Let $X$ be a quasi-compact and quasi-separated scheme. The category of
constructible sheaves of sets is the full subcategory of $\Sh(X_\etale)$
consisting of sheaves $\mathcal{F}$ which are coequalizers
$$
\xymatrix{
\mathcal{F}_1
\ar@<1ex>[r] \ar@<-1ex>[r]
&
\mathcal{F}_0 \ar[r]
&
\mathcal{F}}
$$
such that $\mathcal{F}_i$, $i = 0, 1$ is a finite coproduct of sheaves of
the form $h_U$ with $U$ a quasi-compact and quasi-separated
object of $X_\etale$.
\end{lemma}

\begin{proof}
In the proof of Lemma \href{https://stacks.math.columbia.edu/tag/03SA}{lemma torsion colimit constructible (Tag 03SA)}
we have seen that sheaves of this form are constructible.
For the converse, suppose that for every constructible sheaf of
sets $\mathcal{F}$ we can find a surjection $\mathcal{F}_0 \to \mathcal{F}$
with $\mathcal{F}_0$ as in the lemma. Then we find our surjection
$\mathcal{F}_1 \to \mathcal{F}_0 \times_\mathcal{F} \mathcal{F}_0$
because the latter is constructible by Lemma \href{https://stacks.math.columbia.edu/tag/03RZ}{lemma constructible abelian (Tag 03RZ)}.

\medskip\noindent
By Topology, Lemma
\href{https://stacks.math.columbia.edu/tag/09Y4}{lemma constructible partition refined by stratification (Tag 09Y4)}
we may choose a finite stratification
$X = \coprod_{i \in I} X_i$ such that $\mathcal{F}$ is finite locally
constant on each stratum. We will prove the result by induction on
the cardinality of $I$. Let $i \in I$ be a minimal element in the
partial ordering of $I$. Then $X_i \subset X$ is closed.
By induction, there exist finitely many quasi-compact and quasi-separated
objects $U_\alpha$ of $(X \setminus X_i)_\etale$ and a surjective
map $\coprod h_{U_\alpha} \to \mathcal{F}|_{X \setminus X_i}$.
These determine a map
$$
\coprod h_{U_\alpha} \to \mathcal{F}
$$
which is surjective after restricting to $X \setminus X_i$. By
Lemma \href{https://stacks.math.columbia.edu/tag/03RV}{lemma characterize finite locally constant (Tag 03RV)}
we see that $\mathcal{F}|_{X_i} = h_V$ for some scheme $V$ finite \'etale
over $X_i$. Let $\overline{v}$ be a geometric point of $V$ lying
over $\overline{x} \in X_i$. We may think of $\overline{v}$ as an
element of the stalk $\mathcal{F}_{\overline{x}} = V_{\overline{x}}$.
Thus we can find an \'etale neighbourhood $(U, \overline{u})$
of $\overline{x}$ and a section $s \in \mathcal{F}(U)$ whose stalk at
$\overline{x}$ gives $\overline{v}$. Thinking of $s$ as a map
$s : h_U \to \mathcal{F}$, restricting to $X_i$ we obtain a morphism
$s|_{X_i} : U \times_X X_i \to V$ over $X_i$ which maps $\overline{u}$
to $\overline{v}$. Since $V$ is quasi-compact (finite over the closed
subscheme $X_i$ of the quasi-compact scheme $X$) a finite number
$s^{(1)}, \ldots, s^{(m)}$ of these sections of $\mathcal{F}$ over
$U^{(1)}, \ldots, U^{(m)}$ will determine a jointly
surjective map
$$
\coprod s^{(j)}|_{X_i} : \coprod U^{(j)} \times_X X_i \longrightarrow V
$$
Then we obtain the surjection
$$
\coprod h_{U_\alpha} \amalg \coprod h_{U^{(j)}} \to \mathcal{F}
$$
as desired.
\end{proof}
\end{document}
